% Teil II: der durchgerechnete Weg.

% Gedankenschritt: Python-Wert vs. Formel numerisch, dann Induktion mit
% Nenner 6 und Ausklammern von (n+1).

\subsection*{Gegeben und gesucht}
(a) Programm für $\sum_{k=0}^{n}k^2$; (b) mit $n(n+1)(2n+1)/6$
vergleichen; (c) die Formel beweisen.

\subsection*{a) Programm}
\begin{verbatim}
sum([k**2 for k in range(n+1)])
\end{verbatim}
Für $n=5$: $0+1+4+9+16+25=55$. \checkmark

\subsection*{b) Vergleich}
\begin{center}
\begin{tabular}{c|ccc}
$n$ & $\sum k^2$ & $\frac{n(n+1)(2n+1)}6$ & überein\\
\hline
1 & 1 & $\frac{1\cdot2\cdot3}{6}=1$ & ja\\
2 & 5 & $\frac{2\cdot3\cdot5}{6}=5$ & ja\\
3 & 14 & $\frac{3\cdot4\cdot7}{6}=14$ & ja\\
5 & 55 & $\frac{5\cdot6\cdot11}{6}=55$ & ja
\end{tabular}
\end{center}
Werte stimmen für alle getesteten $n$ überein --- die Vermutung wird
gestützt (bewiesen wird sie durch (c)).

\subsection*{c) Beweis per Induktion}
\emph{IA} ($n=1$): $1=\frac{1\cdot2\cdot3}6$. \checkmark
\emph{IS} ($n\to n+1$), mit der Annahme
$\sum_{k=0}^{n}k^2=\frac{n(n+1)(2n+1)}6$: Der neue Summand ist
$(n+1)^2$; wir schreiben ihn auf den Nenner $6$:
\begin{align}
  \sum_{k=0}^{n+1}k^2
  &=\frac{n(n+1)(2n+1)}6+(n+1)^2 \\
  &=\frac{n(n+1)(2n+1)+6(n+1)^2}{6} \\
  &=\frac{(n+1)\bigl(n(2n+1)+6(n+1)\bigr)}6 \\
  &=\frac{(n+1)(2n^2+7n+6)}6 \\
  &=\frac{(n+1)(n+2)(2n+3)}6 .
\end{align}
Das ist die Behauptung für $n+1$ (man muss $2n^2+7n+6$
faktorisieren: $(n+2)(2n+3)$).

\begin{ergebnis}
\begin{equation*}
  \sum_{k=0}^{n} k^2=\frac{n(n+1)(2n+1)}6\quad\text{für alle }n\ge1 .
\end{equation*}
\end{ergebnis}

\begin{kontrolle}
$n=4$: $0+1+4+9+16=30$, Formel $4\cdot5\cdot9/6=30$; $n=6$:
$0+\dots+36=91$, Formel $6\cdot7\cdot13/6=91$. Passt.

\textbf{Grenzfall/Analyse des Faktors:} Der Zähler $n(n+1)(2n+1)$ ist
stets durch $6$ teilbar (unter drei aufeinanderfolgenden Faktoren
$n(n+1)(n+2)$ gibt es eine gerade Zahl und eine durch $3$; $2n+1$
ist mit $n$ und $n+1$ ``ergänzt`` zu dieser Folge). Dass das Ergebnis
eine ganze Zahl ist, ist ein guter Konsistenztest der Formel.
\end{kontrolle}

\begin{merksatz}
Python prüft die Formel numerisch, die Induktion beweist sie: im IS den
Term $(n+1)^2$ auf den gemeinsamen Nenner $6$ erweitern und $n+1$
ausklammern; das Skript \texttt{sum([k**2 for k in range(n+1)])}
rechnet genau dieselbe Summe wie die Formel.
\end{merksatz}
